%\DeclareGraphicsExtensions{.png}
\chapter{Redox Reactions}
\def\shorttitle{Redox Reactions} 
\section[Background]{Background}

Understanding oxidation-reduction (redox) reactions that occur in natural 
and contaminated systems is central to developing accurate conceptual and 
quantitative models. Redox reactions involve the transfer of electrons 
between constituents and therefore, changes in the oxidation state of 
one or more elements. Changes in oxidation state of an element can 
result in dramatic changes in its chemical and physical properties.  
For example, iron in the plus three oxidation 
state - Fe(III) - in surficial deposits renders 
a reddish hew to Mars whereas iron in the plus II oxidation 
state - Fe(II) - in surficial basalts on the Moon is, in no 
small part, responsible for the blackness of the lunar mares. 
Two important examples of redox reactions are the oxidation 
of aqueous Fe(II) to Fe(III) or organic carbon by dissolved oxygen:

\begin{equation}
4Fe^{+2} + O_{2,aq} + 10H_2O \rightarrow 4Fe(OH)_{3,s} + 8H^+  \\
\label{eq:fe2_oxidation}
\end{equation}

\begin{equation}
CH_2O  + O_{2,aq} + 10H_2O \rightarrow CO_{2,aq} + H_2O \\
\label{eq:om_reduction}
\end{equation}

Rather than representing formaldehyde, $CH_2O$ is often used to represent 
one sixth of a glucose molecule.

The classical approach to discussing redox reactions in aqueous geochemistry 
is to split the redox reactions into two hypothetical 
half reactions involving electrons. 
Using this approach, reaction \ref{eq:fe2_oxidation} and \ref{eq:om_reduction} 
would be broken up into the following half reactions:

\begin{equation}
Fe^{+2} + 3H_2O  \rightarrow Fe(OH)_{3,s} + 3H^+ + e^-  \\
\label{eq:fe_half_reac}
\end{equation}

\begin{equation}
CH_2O + H_2O \rightarrow CO_{2,aq} + 4H^+ + 4e^-  \\                    
\label{eq:om_half_reac}
\end{equation}

\begin{equation}
O_{2,aq} + 4H^+ + 4e^-  \rightarrow 2H_2O \\
\label{eq:ox_half_reac}
\end{equation}

Mass action expressions are then written for the half reactions.
For example, for the Fe(II)/Fe(III) half reaction, 
the mass action expression is:

\begin{equation}
K = a_{Fe(OH)_{3,s}} a_{H^+}^3 a_{e^-}/a_{Fe}^{+2} a_{H_2O}^3 \\
\label{eq:rdx5}
\end{equation}

The standard state of the $e^-$ is the same as other aqueous species,
1 mole per kilogram water.  
This approach comes from the field of electrochemistry, where the oxidation and 
reduction half reactions can be made to occur in separate place
connected, for example, 
through a meter that measures the electromotive force (or power source) 
and a salt bridge.
In dilute solutions the activity of water can be equated with its mole fraction,
which is essentially 1. If pure ferric hydroxide forms then it 
can be assumed to be in 
its standard state (mole fraction $Fe(OH)_{3,s}$
equals 1), 
so \ref{eq:rdx6} can be re-written:

\begin{equation}
K=a_{H^+}^3 a_{e^-}/a_{{Fe}^{+2}} \\    
\label{eq:rdx7}
\end{equation}
                                                         
Taking the base-10 logarithm, \ref{eq:rdx7} becomes: 
\begin{equation}
logK = 3loga_{H^+} + loga_{e^-} -loga_{{Fe}^{+2}}  \\                          
\label{eq:rdx8}
\end{equation}

Analogous to the definition of $pH$,
$loga_{e^-}$
is defined as $pe$.
This formalism is useful for graphically representing the 
relative thermodynamic stability of redox-reactive species 
in natural and contaminated systems 
(typically in their standard state).

Electrochemists have established logK values for 
half reactions based on the 
convention that the hydrogen gas -
hydrogen ion half reaction
\begin{equation}
H_{2,g}  =  2H^+ + 2e^- \\   
\label{eq:rdx9}
\end{equation} 
has an electromotive force and, therefore a logK 
value of 0 for $H_{2,g}$ and $H^+$ in their 
standard states (1 atmosphere and 1 mol/kgw, respectively).  
Using this convention, logK values for half reactions can be assembled 
in order of thermodynamic favorability, 
called a redox ladder, such as Table \ref{tab:rdx1}.  
Note that the logK values apply to all species in their standard state 
(1 mol/kgw for aqueous species, pure solid phase activity 
equals 1, and 1 atmosphere for gasses).  
Also, the logK value for half reactions where solid phases are participants 
depends on the thermodynamic stability of the solid phase.  
For example, reduction of ferrihydrite (fhy) is more favorable than 
reduction of goethite (based on the logK values for goethite that 
applies to the Cape Cod aquifer, Kent et al., 2008) because goethite 
is thermodynamically more stable than ferrihydrite.  
Plots of $pH$ against $pe$ constructed using reactions and logK values like 
those in Table \ref{tab:rdx1} are used to illustrate 
trends expected in redox processes 
with changes in pH. In the exercise that follows, we will explore these trends.

\begin{table}[h]
\topcaption{Half reactions and logK values for some redox-reactive species}
\centering
\footnotesize
%\begin{tabular*}{0.4\textwidth}{@{\extracolsep{\fill}}lll}
\begin{tabular}{cc}
\toprule
Half reactions, per electron																& logK \\ \hline
0.25O2,aq + e- + H+ = 0.5H2O																& 21.52 \\
0.2NO3- + e- + 1.2H+ = 0.1N2,aq + 0.6H2O										& 20.71 \\
0.5MnO2:H2Os,pyrolucite + e- + 2H+ = 0.5Mn+2 + 1.5H2O				& 20.69 \\
0.5HAsO4-2 + e- + 2H+   =  0.5As(OH)3,aq +  0.5H2O					& 19.80 \\
FeOOH + e- + 3H+ = Fe+2 + 2H2O  fhy													& 17.91 \\
FeOOH + e- + 3H+ = Fe+2 + 2H2O  Goethite (Kent et al. 2008)	& 13.52 \\
0.125SO4-2 + e- + 1.125H+ = 0.125HS- + 0.5H2O								& 4.21 \\
0.125CO3-2 + 1.25H+ + e- = 0.125CH4.aq + 0.375H2O						& 5.13 \\
\bottomrule
\label{tab:rdx1}
\end{tabular}
\end{table}	

\section[PHREEQC Exercise 6]{PHREEQC Exercise 6}

This series of exercises involves using PHREEQC to react 
oxidized and reduced species. The oxidation reduction reactions 
effect changes in pH and other chemical conditions.

\begin{table}[h]
\caption{Water compositions for the mixing exercise. Adjust Na to achieve charge balance for the oxic solution and 
Cl to achieve charge balance for the anoxic solution.}
\begin{center}
\begin{tabular}{lccc} 
\hline
Parameter   & Oxic GW & Anoxic GW & Units \\ \hline
pH          & 6.50 & 6.50 & -      \\
pe          & 12   & 4    & -      \\
Temperature & 21   & 21   & C      \\
Alkalinity  & 710  & 710  & ueq/L  \\
Na          & 780  & 580  & umol/L \\ 
K           & 8.0  & 75   & umol/L \\
Mg          & 18.0 & 53   & umol/L \\
Ca          & 9.0  & 170  & umol/L \\ 
Cl          & 82   & 550  & umol/L \\ 
S(6)        & 24.0 & 100  & umol/L \\ 
O(0)        & 500  & 0    & umol/L \\ 
Fe(2)       & 0    & 179  & umol/L \\ \hline
\end{tabular}
\footnote{
\label{redox_phc_table}
\end{center}
\end{table}


\subsection[Mixing-induced redox reactions]{Mixing-induced redox reactions}

\begin{itemize}
\item Set up two solutions, one with $O_2$ and one with $Fe(II)$, using the compositions in Table \ref{redox_phc_table}. {\color{red} \textbf{MIX}} them in equal proportions.  (Keywords: {\color{red} \textbf{SOLUTION n}}, {\color{red} \textbf{SOLUTION m}}, {\color{red} \textbf{MIX o}}.).  
The solutions begin with the same $pH$ and $alkalinity$. What about the $TCO_2$ concentration ? 
$TCO_2$ is the sum of the concentrations of all carbonate-containing aqueous species. 
	\begin{enumerate} 
	\item What happens to the $pH$, $alkalinity$, and $TCO_2$ concentrations after mixing ?
	\item What happens to the concentrations of $O_2$, $Fe(II)$, and $Fe(III)$ after mixing?  What are the dominant $Fe$ aqueous species?
	\item What happens to the $PCO_2$ after mixing?
	\end{enumerate} 
\item Use the anoxic solution. Allow 10 umoles of $O_2(g)$ from the atmosphere ($PO_2$ = 0.2 atmospheres) to equilibrate with the solution. 
(Keyword: {\color{red} \textbf{EQUILIBRIUM\_PHASES n}}.)
  \begin{enumerate} 
	\item What happens to the $pH$, $alkalinity$, and $TCO_2$ concentrations after mixing ?
	\item What happens to the concentrations of $O_2$, $Fe(II)$, and $Fe(III)$ after mixing?  What are the dominant $Fe$ aqueous species?
	\item What happens to the $PCO_2$ after mixing?
  \end{enumerate}
\item Repeat the previous exercise where first 100 $\mu$moles then 1000 $\mu$moles of $O_2$ from the atmosphere equilibrate with the anoxic solution, but allow ferrihydrite (Fe(OH)3(a)) or goethite (Goethite) to precipitate until it reaches equilibrium.
	\begin{enumerate}  
	\item What happens to the $pH$, $alkalinity$, and $TCO_2$ concentrations after mixing ?
	\item What happens to the concentrations of $O_2$, $Fe(II)$, and $Fe(III)$ after mixing?  What are the dominant $Fe$ aqueous species ?
	\item What happens to the $PCO_2$ after mixing ?
	\item What happens when the solution is allowed to equilibrate with atmospheric $PCO_2$ ?
	\end{enumerate} 
\item Simulate reacting the oxic solution with 500 umoles of amorphous $FeS$ ($FeS(ppt)$) (~44 mg) allowing $Fe(OH)3(a)$ to precipitate. 
	\begin{enumerate}  
	\item What got oxidized ? What did not get oxidized? What was the $pH$ ?
	\item Repeat the simulation, this time allowing the suspension to equilibrate with atmospheric $O_2$. What was the $pH$ ? What got oxidized ?
	\end{enumerate}  
\end{itemize} 

\subsection[Redox titration]{Redox titration}

The next PHREEQC exercise compares the standard sequence of oxidation-reduction reactions 
predicted using the redox ladder to the sequence in which oxidants (electron acceptors) 
actually get reduced by organic carbon based on thermodynamic considerations.   
The exercise uses the keyword {\color{red} \textbf{REACTION n}} to add $CH_2O$, 
which represents one sixth of a glucose molecule, to a suspension with aqueous and solid-phase oxidants.  
As $CH_2O$ is added it reacts with the oxidants that are most thermodynamically favorable. T
hermodynamic favorability depends on concentrations and solid-phase solubilities 
(which depend on $pH$ for some important oxidants).  
Thus the order in which the oxidants react with $CH_2O$ may change during the course 
of the redox titration as the chemical conditions change.  

\textbf{Note, that you will need to use either the {\color{red} \textbf{Wateq4f.dat}} 
or {\color{red} \textbf{Minteq.v4.dat}} database, 
because the Phreeqc.dat database does not include arsenic.}


\begin{itemize}
\item Create a PHREEQC input file and set up the initial solution and solid-phase compositions 
in the suspension using the data in Table 6.3.  
Assume the suspension has one kilogram of water.  
(Keywords: {\color{red} \textbf{SOLUTION n}}, {\color{red} \textbf{EQUILIBRIUM_PHASES n}}).  
Adjust the chloride concentration to insure charge balance of the solution. 
Two solid phases with reduced species may precipitate if they become supersaturated: 
amorphous FeS (FeS(ppt) in Phreeqc) and siderite.
\end{itemize}















